在 \(P\) 本身的极化缺陷之外，\(V_3'\)-等型块 \(A_{3'}\sim_{\QQ} P^3\) 的主极化化代价、全部额外有理端同态以及正维子块类型也可由同一分解立即锁定。

\begin{theorem}[Prym 三重等型块的主极化化、条件性矩阵端代数与子块分类]\label{thm:xi-time-part9n-prym-cubed-principalization-endomorphism}
\leanverified{paper\_xi\_time\_part9n\_prym\_cubed\_principalization\_endomorphism}
沿用结论 \ref{con:xi-time-part9n-prym-polarization-minimal-isogeny-degree} 的记号，记
\[
P:=\operatorname{Prym}(Y/Z),
\qquad
A_{3'}\sim_{\QQ} P^3,
\qquad
\mathrm{type}(\lambda_P)=(1^4,2^5).
\]
令 \(\lambda_{P^3}\) 表示 \(P^3\) 上的乘积极化。则成立：
\begin{enumerate}
  \item \(\lambda_{P^3}\) 的极化型为
  \[
  \mathrm{type}(\lambda_{P^3})=(1^{12},2^{15}).
  \]
  \item 若 \((B,\Theta)\) 为主极化阿贝尔簇，且
  \[
  f:P^3\to B
  \]
  为同源并满足
  \[
  f^*\Theta=\lambda_{P^3},
  \]
  则
  \[
  \deg f=2^{15}.
  \]
  特别地，\(P^3\) 的主极化化最小次数恰为 \(2^{15}\)。
  \item 若猜想 \ref{conj:killo-minimal-extra-endomorphism-large-prym} 的第一条成立，即 \(P\) 在 \(\overline{\QQ}\) 上绝对单，且
  \[
  \operatorname{End}_{\overline{\QQ}}(P)=\ZZ,
  \]
  则
  \[
  \operatorname{End}^0_{\QQ}(A_{3'})\cong M_3(\QQ),
  \qquad
  Z\!\bigl(\operatorname{End}^0_{\QQ}(A_{3'})\bigr)=\QQ.
  \]
  并且 \(A_{3'}\) 的任意正维 \(\QQ\)-阿贝尔子簇都同源于
  \[
  P^r,\qquad 1\le r\le 3.
  \]
  特别地，\(A_{3'}\) 的任意 \(9\) 维简单 \(\QQ\)-因子都唯一同源于 \(P\)。
\end{enumerate}
\end{theorem}
\begin{proof}
乘积极化的 elementary divisors 由各因子的 elementary divisors 直接并列给出，故
\[
\mathrm{type}(\lambda_{P^3})
=
(1^4,2^5)^{\times 3}
=(1^{12},2^{15}).
\]
这证明了第 \(1\) 条。

由极化型公式，
\[
\deg(\lambda_{P^3})=(2^{15})^2.
\]
另一方面，若 \(\Theta\) 为主极化且 \(f^*\Theta=\lambda_{P^3}\)，则
\[
\deg(\lambda_{P^3})
=
\deg(f^*\Theta)
=
(\deg f)^2\deg(\Theta)
=
(\deg f)^2.
\]
故
\[
\deg f=2^{15}.
\]
再者，\(\ker(\lambda_{P^3})\) 作为有限辛群含有阶为 \(2^{15}\) 的极大各向同性子群 \(K\)。对商映射
\[
q:P^3\to P^3/K
\]
应用极化下降的标准构造，可得某个主极化 \(\Theta_K\) 满足
\[
q^*\Theta_K=\lambda_{P^3}.
\]
因此 \(P^3\) 的主极化化确可由次数 \(2^{15}\) 的同源实现，从而该次数既是必要下界，也是可达下界，第 \(2\) 条成立。

若猜想 \ref{conj:killo-minimal-extra-endomorphism-large-prym} 的第一条成立，则
\[
\operatorname{End}^0_{\overline{\QQ}}(P)
=
\operatorname{End}_{\overline{\QQ}}(P)\otimes_{\ZZ}\QQ
=
\QQ.
\]
因而
\[
\operatorname{End}^0_{\QQ}(P)=\QQ.
\]
另一方面，由结论 \ref{con:xi-terminal-zm-s4-jacobian-a3p-isog-prym3-cubed}，
\[
A_{3'}\sim_{\QQ} P^3.
\]
同源阿贝尔簇的有理端同态代数在 \(\QQ\)-同源范畴中同构，而任意阿贝尔簇 \(A\) 都满足
\[
\operatorname{End}^0_{\QQ}(A^3)\cong
M_3\!\bigl(\operatorname{End}^0_{\QQ}(A)\bigr).
\]
取 \(A=P\) 即得
\[
\operatorname{End}^0_{\QQ}(A_{3'})
\cong
\operatorname{End}^0_{\QQ}(P^3)
\cong
M_3(\QQ).
\]
其中心显然等于 \(\QQ\)。

再设 \(C\subseteq A_{3'}\) 为任意正维 \(\QQ\)-阿贝尔子簇。经由 \(A_{3'}\sim_{\QQ} P^3\) 的同源识别，\(C\) 在 \(\QQ\)-同源范畴中对应于某个幂等元
\[
e\in \operatorname{End}^0_{\QQ}(P^3)\cong M_3(\QQ)
\]
的像。线性代数表明，\(M_3(\QQ)\) 中任一幂等元都按共轭同价于
\[
\operatorname{diag}(I_r,0)\qquad (0\le r\le 3),
\]
其中 \(r=\operatorname{rank}(e)\) 唯一确定。于是
\[
\operatorname{Im}(e)\sim_{\QQ} P^r.
\]
由于 \(C\) 正维，故 \(1\le r\le 3\)。这就给出了全部正维 \(\QQ\)-阿贝尔子簇的同源类型分类。

最后，若 \(C\) 还是简单且 \(\dim C=9=\dim P\)，则只能有 \(r=1\)，从而
\[
C\sim_{\QQ} P.
\]
故第 \(3\) 条全部成立。证毕。
\end{proof}

\endinput
